\documentclass[11pt,a4paper]{article}
\usepackage[utf8]{utf8}
\usepackage[spanish]{babel}
\usepackage{graphicx}
\usepackage{booktabs}
\usepackage{geometry}
\geometry{margin=2.5cm}

\title{\textbf{Proyecto Golden Hour: Reporte de Accesibilidad Médica}}
\author{Lambayeque, Ayacucho y San Martín}
\date{\today}

\begin{document}
\maketitle

\section{Introducción}
Este informe presenta el análisis geoespacial de tiempos de traslados viales hacia establecimientos de salud de nivel resolutivo (II-1 a III-E) para los departamentos de Lambayeque, Ayacucho y San Martín.

\section{Resultados}
La Figura \ref{fig:dist} exhibe la distribución de tiempos de traslados, mientras que la Tabla \ref{tab:worst} resume los distritos con mayores limitaciones de infraestructura.

\begin{figure}[htbp]
    \centering
    \includegraphics[width=0.85\linewidth]{figures/fig_access_distribution.pdf}
    \caption{Distribución distrital de accesibilidad médica.}
    \label{fig:dist}
\end{figure}

\begin{table}[htbp]
    \centering
    \caption{Distritos prioritarios según brechas de acceso}
    \input{../data/outputs/table_worst_districts.tex}
    \label{tab:worst}
\end{table}

\end{document}

